\documentclass[10pt,a4paper]{amsart}
\usepackage[margin=19mm]{geometry}
\usepackage{amsmath,amssymb,mathtools}
\usepackage[scr=rsfs,cal=euler]{mathalfa}
\usepackage{xfrac}
\usepackage[unicode,hidelinks,bookmarks=false]{hyperref}
\newcommand{\ie}{i.e.\ }
\newcommand{\cf}{cf.\ }
\newcommand{\resp}{respectively\ }
\newcommand{\Subsection}{\subsection}
\providecommand{\nobreakdash}{\nobreak\hbox{-}}
\setlength{\emergencystretch}{2em}
\multlinegap0pt
\usepackage{amssymb}
\usepackage{mathtools}
\usepackage{bm}
\usepackage{tikz}
\newtheorem{proposition}{Proposition}
\newcommand{\Aset}{\mathcal{A}}
\newcommand{\R}{\mathbb{R}}
\newcommand{\Sset}{\mathcal{S}}
\DeclareMathOperator*{\argmax}{arg\,max}
\title{Decision Making Needs Uncertainty Quantification [Lecture Notes]}
\author{Osvaldo Simeone}
\date{Source: arXiv:2607.14407v4 (CC BY 4.0)}
\begin{document}
\maketitle

\noindent\emph{Source and licence.} Decision Making Needs Uncertainty Quantification [Lecture Notes]. Version arXiv:2607.14407v4 (CC BY 4.0), distributed by arXiv under the Creative Commons Attribution 4.0 International licence, http://creativecommons.org/licenses/by/4.0/, as recorded in the arXiv licence metadata for this exact version and shown on the abstract page https://arxiv.org/abs/2607.14407v4. Full licence notice: \texttt{licenses/arxiv-2607.14407-CC-BY-4.0.txt}.\par\smallskip

\noindent\textit{Editorial scope.} Counted unit: the article's proposition prop:ra (source0.tex lines 371--384). The article's complete original proof of it is delivered with it (source0.tex lines 385--406, one \texttt{proof} environment of 22 lines). No mathematics is written, repaired or re-derived: the statement and the proof are the article's own lines, in the article's own notation, and no retained line is substituted or re-flowed.  Retained uncounted as the source-local context the proof needs: lines 194--196; lines 283--287; lines 291--293; lines 303--306; lines 309--322.  One declared adaptation: exactly 1 ancillary strings inside the retained spans are omitted (image-inclusion commands and, where the source repeats a label, the repeated label definitions) and no other line differs from the source; the figure environments, with their placement, captions and labels, are kept, and the package records the omitted strings in fidelity receipts bound to the affected bodies.  No retained line contains a citation, so the wrapper carries no bibliography and no bibliography entry.  The retained lines contain the article's own diagram code, which the wrapper typesets with the standard tikz packages; no image file is delivered and the package declares no asset dependency.  Editorial formatting only, in addition: an AMS \texttt{amsart} preamble that restates the article's own declarations and macros used by the retained lines, namely the environments \texttt{proposition} on separate counters, together with the standard packages those lines need.  The retained environments print in this document as Proposition 1 in order of appearance, whereas the article prints the same items with the numbering of its own full document.  The complete native source of the article as distributed in the arXiv e-print for this exact version is bundled unchanged as \texttt{sources/205292/source0.tex}.
\begin{equation}\label{eq:optaction}
a^*(o)=\argmax_{a \in \Aset} U(a| o).
\end{equation} An example of utility $u(s,a)$, posterior $p(s|o)$, and average utility $U(a|o)$ is shown in Fig. \ref{fig:ex}.

\begin{equation}
\label{eq:var}
V_\alpha(a| o)=\sup\Bigl\{\nu\in\R:\
\Pr_{s\sim p(s| o)}\!\left[\,u(s,a)\ge\nu\,\right]\ge 1-\alpha\Bigr\}.
\end{equation}

\begin{equation}\label{eq:optactionalpha}
a_\alpha^*(o)=\argmax_{a \in \Aset} V_\alpha(a| o).
\end{equation}

\begin{equation}
\label{eq:coverage}
\Pr_{s\sim p(s| o)}\!\left[\,s\in \mathcal{C}_{\alpha}(o)\,\right]\ge 1-\alpha.
\end{equation} By (\ref{eq:coverage}), the set $\mathcal{C}_{\alpha}(o)$ is guaranteed to contain the true state $s$ with probability no smaller than $1-\alpha$ given the observation $o$. An example of a level-$\alpha$ prediction set is shown in Fig. \ref{fig:ex}.

\begin{figure}[t]
  \centering
    % or a PDF (risk_decision.pdf)
  \caption{Risk-neutral versus risk-averse decision making for a fixed
  observation~$o$. The  blue curves are level sets of the utility
  $u(s,a)$ over the state $s$ (vertical axis) and the action $a$
  (horizontal axis). The posterior $p(s|o)$ is shown against the
  state axis (left); the level-$\alpha$ prediction set $C_\alpha(o)$
  collects states carrying posterior mass $1-\alpha$ (hatched). 
  The average utility $U(a| o)$
  (solid) and the value-at-risk $V_\alpha(a| o)$ (dashed) yield the optimal action $a^*(o)$ in~\eqref{eq:optaction} and 
  $a^*_\alpha(o)$ in \eqref{eq:optactionalpha}, respectively. }
  \label{fig:ex}
\end{figure}

\begin{proposition}[Prediction sets are the optimal interface for risk-averse decisions]
\label{prop:ra}
The optimal risk-averse action \eqref{eq:optactionalpha} is the max--min
action of an optimally chosen level-$\alpha$ prediction set, i.e., 
\begin{equation}
\label{eq:raset}
a_\alpha^*(o)=a_{\mathrm{mm}}\!\left(C_\alpha^*(o)\right),
\end{equation}
where the optimized prediction set is given by
\begin{equation}
\label{eq:optset}
C_\alpha^*(o)\in\!\!\argmax_{\substack{\mathcal{C}\subseteq\Sset:\\ \Pr_{s\sim p(s| o)}[s\in \mathcal{C}]\ge 1-\alpha}}\!\!U_{\mathrm{mm}}(\mathcal{C}).
\end{equation}
\end{proposition}

\begin{proof}
We first establish, for a fixed action $a$, the identity
\begin{equation}
\label{eq:varset}
V_\alpha(a| o)=\!\!\max_{\substack{\mathcal{C}\subseteq\Sset:\\ \Pr_{s\sim p(s| o)}[s\in \mathcal{C}]\ge 1-\alpha}}\!\!\min_{s\in \mathcal{C}}u(s,a),
\end{equation}
which states that the value-at-risk equals the best worst-case utility achievable over a
covering set. Given this result, optimizing the value-at-risk over the action $a$ yields the max-min action $a_{\mathrm{mm}}\!\left(C_\alpha^*(o)\right)$, as claimed in the proposition. 

For the ``$\ge$'' direction in \eqref{eq:varset}, take any
feasible set $\mathcal{C}$, satisfying the coverage condition (\ref{eq:coverage}), and let $m(\mathcal C)=\min_{s\in \mathcal{C}}u(s,a)$. Then, we have the inclusion relationship 
$\{s\in \mathcal{C}\}\subseteq\{u(s,a)\ge m(\mathcal C)\}$, which implies the inequality 
$\Pr_{s\sim p(s| o)}[u(s,a)\ge m(\mathcal C)]\ge\Pr_{s\sim p(s| o)}[s\in \mathcal{C}]\ge1-\alpha$. In turn,  by the
definition \eqref{eq:var} of the value-at-risk, this gives the inequality 
$m(\mathcal{C})\le V_\alpha(a| o)$, which holds also when optimizing over the set. This proves the desired  $\ge$ inequality in \eqref{eq:varset}.

For the reverse direction, consider the set
$\mathcal{C}_0=\{s:u(s,a)\ge V_\alpha(a|o)\}$. By definition of value-at-risk, the set $\mathcal{C}_0$ is  feasible for problem \eqref{eq:varset}, and
satisfies the inequality $\min_{s\in C_0}u(s,a)\ge V_\alpha(a|o)$, proving the desired $\le$ inequality in \eqref{eq:varset}.


\end{proof}

\end{document}
